\chapter{Schede di ripasso}
\label{cap:schede}

% ---------------------------------------------------------------------
\section{Severità di un data breach (metodologia ENISA 2013)}

\[
SE = DPC \times EI + CB
\]

\begin{center}
\small
\begin{tabularx}{\linewidth}{l X}
\toprule
\textbf{Fattore} & \textbf{Valori} \\ \midrule
\textbf{DPC} -- Data Processing Context & 1 semplici $\cdot$ 2 comportamentali $\cdot$ 3 finanziari
  $\cdot$ 4 sensibili; aumentabile o riducibile in base al contesto (minori, sicurezza, volume),
  motivando la scelta \\
\textbf{EI} -- Ease of Identification & 0,25 trascurabile $\cdot$ 0,5 limitata $\cdot$
  0,75 significativa $\cdot$ 1 massima \\
\textbf{CB} -- Circumstances of the Breach & riservatezza 0--0,5 $\cdot$ integrità 0--0,5 $\cdot$
  disponibilità 0--0,5 $\cdot$ intento malevolo +0,5 \\
\midrule
\textbf{Scala} & $SE<2$ bassa $\cdot$ $2 \le SE < 3$ media $\cdot$ $3 \le SE < 4$ alta $\cdot$
  $SE \ge 4$ molto alta \\ \bottomrule
\end{tabularx}
\end{center}

\begin{esame}
\begin{itemize}
  \item Non conta ``indovinare il numero'' ma \textbf{motivare} ogni valore (accountability).
  \item La metodologia non distingue minori e adulti: se ci sono soggetti vulnerabili (minori,
    pazienti, lavoratori), conviene alzare il DPC motivando e comunicare comunque agli interessati.
  \item Per i minori la comunicazione ex art.~34 va ai \textbf{genitori}.
  \item Anche l'indisponibilità temporanea dei dati è un data breach.
\end{itemize}
\end{esame}

\begin{nota}
Scala tratta dalle slide del caso Pulse K (A.A.~2026/27). Gli appunti dell'anno precedente usano
soglie leggermente diverse (es.\ 2--2,5 media, 2,5--3 medio-alta): fa fede la versione di quest'anno.
\end{nota}

% ---------------------------------------------------------------------
\section{Ruoli nelle diverse norme}

\begin{center}
\small
\begin{tabularx}{\linewidth}{l X}
\toprule
\textbf{Norma} & \textbf{Ruoli principali} \\ \midrule
GDPR    & Titolare (decide finalità e mezzi; notifica i breach) $\cdot$ responsabile (tratta per conto
  del titolare; avvisa il titolare senza ingiustificato ritardo) $\cdot$ sub-responsabile $\cdot$
  contitolari $\cdot$ interessato $\cdot$ DPO \\
NIS2    & Soggetti essenziali e importanti $\cdot$ organi di gestione (responsabilità non delegabile)
  $\cdot$ punto di contatto NIS $\cdot$ CSIRT e autorità competente \\
DORA    & Entità finanziarie $\cdot$ fornitori terzi di servizi ICT $\cdot$ organo di gestione
  $\cdot$ funzione di controllo del rischio ICT \\
CRA     & Fabbricante $\cdot$ importatore $\cdot$ distributore \\
AI Act  & Provider $\cdot$ deployer $\cdot$ importatore $\cdot$ distributore $\cdot$ fabbricante del
  prodotto (art.~25: riqualificazione come provider) \\ \bottomrule
\end{tabularx}
\end{center}

% ---------------------------------------------------------------------
\section{Date da ricordare}

\begin{center}
\small
\begin{tabularx}{\linewidth}{l X}
\toprule
\textbf{Data} & \textbf{Evento} \\ \midrule
25/05/2018 & Il GDPR diventa applicabile \\
17/01/2025 & DORA diventa applicabile \\
11/09/2026 & CRA: si applicano gli obblighi di segnalazione (art.~14) \\
11/12/2027 & CRA: applicazione integrale (requisiti dell'Allegato I) \\
02/12/2027 & AI Act: obblighi per i sistemi ad alto rischio dell'Allegato III (dopo l'AI Omnibus) \\
02/08/2028 & AI Act: obblighi per i sistemi ad alto rischio dell'Allegato I (dopo l'AI Omnibus) \\ \bottomrule
\end{tabularx}
\end{center}

\begin{integrazione}[Altre date]
NIS2: D.lgs.~138/2024 in vigore dal 16/10/2024 (recepimento italiano). AI Act: in vigore dal 01/08/2024;
divieti dell'art.~5 dal 02/02/2025; obblighi GPAI dal 02/08/2025. CSA: in vigore dal 27/06/2019.
Le date dell'AI Omnibus sono dalle slide introduttive (Capitolo~\ref{cap:aiact}).
\end{integrazione}

\dacompletare{Rivedere le date NIS2, DORA e AI Act dopo le relative lezioni.}
